\documentclass[output=paper]{langscibook}
\ChapterDOI{10.5281/zenodo.22880252}
\title{\textit{Karaa} – Pandanus}
\author{Ruth Sima Rigamu}
\abstract{}
\begin{document}
\maketitle\label{chap:13}
\begin{figure}[H]
    \captionsetup[subfigure]{oneside,margin={0cm,0cm}}
    \captionsetup{justification=centering}
    \subfloat[]{\includegraphics[trim=0 0 0 11,clip,width=.35\linewidth]{figures/13_1.jpg}\label{fig:13_1}}
    \subfloat[]{\includegraphics[trim=0 18 0 35,clip,width=.30\linewidth]{figures/13_2.jpg}\label{fig:13_2}}
    \subfloat[]{\includegraphics[trim=30 0 0 0,clip,width=.35\linewidth]{figures/13_3.jpg}\label{fig:13_3}}
    \caption{(a) \textit{Karaa} -- Pandanus, (b) \textit{Varoava} -- Letting dry in the sun,\\ (c) \textit{Vin vonaha} -- Weaving a mat}
\end{figure}
\begin{xltabular}{\textwidth}{lXX}
    (\ref{ex:text9-1})&\textit{\textbf{Karaa}}&\textbf{Pandanus}\\
    (\ref{ex:text9-2})&\textit{O karaa to vin vonaha ni ria ta Teapu.}&The pandanus is used by the Teop people for weaving mats.\\
    (\ref{ex:text9-3})&\textit{Na kosi riori ore paa kahu bono sinoo vaa toon no paka karaa, ore paa vaameapa bari tea roava.}&They cut the pandanus leaves and then scrape the midrib on the back of the pandanus leaves and let them become flexible in the sun.\\
    (\ref{ex:text9-4})&\textit{Me meapa, ore paa rava bari ge eori repaa nahu bari bona tahii.}&And (when the pandanus leaves are) flexible, they heat them over the fire, or they cook them in saltwater.\\
    (\ref{ex:text9-5})&\textit{Murinae eori repaa varoava bari.}&After that they put them in the sun to dry.\\
    (\ref{ex:text9-6})&\textit{Varoava riori, antee beo karaa gaoho komana.}&They put them in the sun until the pandanus is bleached.\\
\end{xltabular}
\begin{figure}[H]
    \captionsetup[subfigure]{oneside,margin={0cm,0cm}}
    \captionsetup{justification=centering}
    \subfloat[]{\includegraphics[trim=0 0 0 375,clip,width=.38\linewidth]{figures/13_4.jpg}\label{fig:13_4}}
    \subfloat[]{\includegraphics[trim=0 18 0 35,clip,width=.62\linewidth]{figures/13_5.jpg}\label{fig:13_5}}
    \caption{(a) \textit{Karaa toutou} -- Thorny pandanus,\\ (b) \textit{Paka toutou} -- Thorny leaves}
\end{figure}
\begin{xltabular}{\textwidth}{lXX}
    (\ref{ex:text9-7})&\textit{\textbf{Karaa Toutou}}&\textbf{Thorny Pandanus} (\textit{Pandanus spiralis})\\
    (\ref{ex:text9-8})&\textit{O karaa toutou na bero kurus nana kasuana teara.}&The thorny pandanus is very common in our coastal area.\\
    (\ref{ex:text9-9})&\textit{Pura pinopino batana mee,}&It grows everywhere without any care,\\
    (\ref{ex:text9-10})&\textit{ei kou o vua nae to oro binubinu batana tahii.}&because its seeds float around everywhere in the sea.\\
    (\ref{ex:text9-11})&\textit{O karaa bona na vaapeha mi nana bono karaa vai o hiki toutou.}&This pandanus is similar to the pandanus without thorns\\
    (\ref{ex:text9-12})&\textit{A taba to vaavaha koa ni nana teo buo karaa bona, o meho na tei mi nana bono toutou teo buo vavaarihi bara toon nae, o meho me ahiki.}&The only thing that is different with these two kinds of pandanus is that one has thorns on the two sides of the leaf and on its back, but the other one does not.\\
    (\ref{ex:text9-13})&\textit{Eara na kahu rara o buo vavaarihi no paka karaa bara toon nae.}&We scrape the two sides of the leaf and its back.\\
    (\ref{ex:text9-14})&\textit{O karaa bona mee o gagaoho pete kurus. O karaa bona o mataa.}&This pandanus is also very white. This pandanus is good.\\
    (\ref{ex:text9-15})&\textit{Na vin hoi bara vonaha ni ria maa moon.}&The women use it for weaving baskets and mats.\\
\end{xltabular}
\begin{figure}[H]
    \includegraphics[trim=0 0 0 0,clip,width=.6\linewidth]{figures/13_6.jpg}
    \caption{\textit{Karaa vaa Solomon} -- Pandanus from the Solomon Islands}
    \label{fig:13_6}
\end{figure}
\begin{xltabular}{\textwidth}{lXX}
    (\ref{ex:text9-16})&\textit{\textbf{Karaa vaa Solomon}}&\textbf{Pandanus from the Solomon Islands}\\
    (\ref{ex:text9-17})&\textit{O karaa vai o hiki toutou na deedee maa rohori Solomon, meara na dao koa rara e bona karaa vaa Solomon.}&The pandanus without thorns was brought from the Solomon Islands, and we just call it the pandanus from the Solomons.\\
    (\ref{ex:text9-18})&\textit{Ahiki to toutou teve.}&It does not have any thorn.\\
    (\ref{ex:text9-19})&\textit{Eara na kahu vatamee koara a toon na maa paka karaa, ore tamee tea rigunu ei beara tau nahue.}&We only scrape the back of the pandanus leaves to make them flexible so that they are flexible to wind them up when we are about to cook them.\\
    (\ref{ex:text9-20})&\textit{O karaa bona o gaohogaoho kurus, o tamee kurus mee.}&This pandanus is very white, but also very flexible.\\
    (\ref{ex:text9-21})&\textit{Na mataa kurus nana tea vin vonaha, hoi ae haata.}&It is very good for weaving mats, baskets and hats.\\
\end{xltabular}
\end{document}